\section{Discussion and Conclusion}
\label{sec:conclusion}

\para{Scope, boundaries, and future work.}
Our findings are grounded in two real-silicon targets on an Intel Sapphire Rapids socket; two hardware boundaries and three open directions merit attention when generalizing:
\begin{enumerate}
  \item \textbf{Host CPU and memory channels:} The exact 330--430 \texttt{CHA WCiLF} cliff (hottest slice $\sim$22--23 entries) is a property of SPR's 32-\texttt{CHA} \texttt{TOR} pool; sockets with different \texttt{TOR} depths (\eg{} Emerald/Granite Rapids or AMD EPYC) will exhibit a different numerical threshold governed by the same shared-request-table exhaustion mechanism. Likewise, our host ran two populated DDR5 channels ($\sim$60\,GB/s); on an 8-channel ($\sim$240\,GB/s) server, the relative STREAM penalty of host-memory DMA rings will shrink proportionally, whereas \texttt{CHA TOR} slot exhaustion under unpaced MMIO push remains independent of DRAM channel count.
  \item \textbf{Device L3 Point-of-Coherence semantics:} Zero-DRAM-penalty \texttt{MT\_NORMAL\_NC} (\texttt{0x44}) relies on CN106's 48\,MiB \texttt{TAD} acting as the CHI Point of Coherence for non-cacheable loads. On ARMv8 SoCs where \texttt{MT\_NORMAL\_NC} bypasses L3 straight to DRAM, designers should instead use our portable coalesced decoupled doorbell (\S\ref{sec:design_rx}) or verify custom \texttt{Inner-NC, Outer-WB} mappings.
  \item \textbf{Future work:} Three genuinely unmeasured directions remain open: (i)~mapping the skewed-burst $M \times K$ trade-off curve on statically partitioned portals (varying active bursting cores $M \in [1, 32]$ and slice width $K \in \{2, 4, 6\}$ under locked \texttt{performance} governors); (ii)~measuring the \texttt{TOR}-exhaustion cliff across additional server microarchitectures (AMD EPYC, ARM Neoverse hosts, and next-generation Xeon); and (iii)~evaluating how CXL~3.0 \texttt{HDM-DB} back-invalidation interacts with paced direct-push rings.
\end{enumerate}

\para{Conclusion.}
MMIO direct push delivers compelling isolation and latency wins---175\,Mpps from a single core with $<$1\% host neighbor interference, sub-microsecond RTT, and +22.5\% to +55\% end-to-end throughput gains on real 100\,GbE key-value services---but only when paired with pacing. Without pacing, open-loop posted writes turn both host over-injection and device frontier collisions into the same backpressure cascade across PCIe or on-die \texttt{IIO}, filling the host's shared \texttt{CHA TOR} past its 330--430-entry cliff and collapsing co-located compute by up to 96.7\%. Combining spatial \texttt{(Physical Core WCB, 64\,B)} line-budgeting on stateless portals with explicit rate, micro-burst, and credit-window pacing on stateful rings---paired with single-binary zero-lookahead \texttt{MT\_NORMAL\_NC} (\texttt{u1\_ldpq}) or coalesced doorbells---supplies the missing piece that makes MMIO direct push fast, predictable, and safe alongside latency-sensitive workloads.
